\section{补充方法与结果}
\label{app:supplement}

\subsection{输入登记与可复现性}
\label{app:inputs}

全部报告结果均由发布清单记录 SHA-256 的 NHSN 原始州周表重建。面板将缺失值保持为缺失，并要求四个连续周日期，不插入零住院。主样本与全部固定窗口由发布版中的确定性键标识。所有保留的表格、图件和正文数值宏均由同一版本化结果包生成；清单记录输入哈希、生成器哈希、实际自助预算、拟合失败与蒙特卡洛诊断。续跑必须匹配完整的源码、输入与配置指纹，否则重新运行。

表~\ref{tab:pooled} 报告合并仿射工作拟合及三种代际间隔变换。3.6 天与固定变异系数设定均明确列为选择；3.2 天、标准差 2.1 天设定复现 \citet{chan2025estimating} 的内禀间隔估计。表~\ref{tab:cv} 报告仅用训练折中心和拟合的逐折留出评分。季节折使用与拟合相同的数据，属描述性结果。

校准登记表将每个输入标明为有据外部估计、由观测面板导出或分析假设。延迟、后代离散度、住院概率与周冲击方差在模拟器中是假设，其作用是结构敏感性而非测量。观察模型以二项住院分配为基准；$\varphi>1$ 表示直接从潜在强度生成、条件方差为 $\varphi\lambda$ 的伽马--泊松计数。模拟器检查条件观察矩与直接计数构造；这些检验验证实现，不证明经验适用性。

\subsection{条件校准}
\label{app:calibration}

每个声明情景在四个固定观测窗口长度和观测首周州计数上模拟。初始零计数保持为零，不抬高；这可能使保留样本小于实证面板，差异在响应诊断中保留。周再生数冲击区分全国与州分量；报告持续性属于另一机制。再生数新息在声明的正下界处截断，并记录截断比例。

对每个实证统计量，用加权直线联系模拟平均截距与声明周方差网格。配对经验抽样与配对模拟面板共同形成条件反演抽样。拟合失败、非有限值、非正隐含值与网格外推均计数并报告；百分位区间只针对有限抽样，并标注条件性。多个情景可给出相近实证摘要而观察机制不同。因此校准报告截距、斜率、样本量、有符号偏差方差比及州内与全国滞后一阶自相关的响应检查，不指定唯一生物学方差或某机制正确的概率。

\subsection{次要敏感性与设计示例}
\label{app:secondary}

表~\ref{tab:secondary} 报告分季截距，以及将各统计量计数除以机构参与率、同时保留原尺度的配对调整。该调整只测记录到的参与率代理；报告机构内部的检测强度、编码与就诊行为未观测。三个夏季摘要未经调整，仅描述单次自愿报告期，不是因果比较或独立重复。本文不声称等效性或报告制度效应。

表~\ref{tab:design} 报告模型条件下的曲线辨别示例。观测季节设计按声明计数倍率循环，并叠加均值为一的交叉对数正态变异；生成曲线为拟合的仿射与幂律均值。校准分位与评价模拟相互独立；在解释功效前先报告零假设拒绝比例及精确二项区间。所有单元的功效都接近 1。这种饱和说明合成生成机制远比真实数据有利，而真实数据中两条曲线无法区分；因此该练习不能回答现实需要多少季节，仅用于记录这一局限。附录~\ref{app:allocation} 的预算表达式保持符号形式；不将周住院比转换为代际设计门槛。

\begin{table}[htbp]
\centering
\small
\caption{合并仿射损失拟合与代际间隔敏感性}\label{tab:pooled}
\setlength{\tabcolsep}{2pt}
\begin{tabularx}{\linewidth}{@{}>{\raggedright\arraybackslash}p{0.2400\linewidth}>{\raggedright\arraybackslash}p{0.1750\linewidth}>{\raggedright\arraybackslash}p{0.1750\linewidth}>{\raggedright\arraybackslash}p{0.1750\linewidth}>{\raggedright\arraybackslash}p{0.1750\linewidth}@{}}
\toprule
统计量 & 参数 & n & 估计值 [95\% CI] & 有限/B \\
\midrule
RA & a & 1748 & 0.01495 [0.01001, 0.02092] & 2000/2000 \\
RA & b & 1748 & 0.56 [0.369, 0.7692] & 2000/2000 \\
RA & b/a & 1748 & 37.44 [18.93, 69.92] & 2000/2000 \\
G & a & 1756 & 0.08401 [0.04986, 0.1234] & 2000/2000 \\
G & b & 1756 & 1.971 [0.9889, 3.042] & 2000/2000 \\
G & b/a & 1756 & 23.46 [8.279, 54.7] & 2000/2000 \\
RA32 & a & 1748 & 0.01191 [0.007997, 0.01668] & 2000/2000 \\
RA32 & b & 1748 & 0.4482 [0.2975, 0.6147] & 2000/2000 \\
RA32 & b/a & 1748 & 37.62 [19.13, 69.89] & 2000/2000 \\
RA32source & a & 1748 & 0.01106 [0.007501, 0.01553] & 2000/2000 \\
RA32source & b & 1748 & 0.4196 [0.2827, 0.5688] & 2000/2000 \\
RA32source & b/a & 1748 & 37.93 [19.42, 69.26] & 2000/2000 \\
\bottomrule
\end{tabularx}
\par\smallskip
\begin{minipage}{\linewidth}\footnotesize\raggedright
有符号百分位区间；比值只在分母为正时定义。逐项保留有限重抽样次数。 RA a: 2000/2000 有限重抽样; RA b: 2000/2000 有限重抽样; RA ratio: 2000/2000 有限重抽样; G a: 2000/2000 有限重抽样; G b: 2000/2000 有限重抽样; G ratio: 2000/2000 有限重抽样; RA32 a: 2000/2000 有限重抽样; RA32 b: 2000/2000 有限重抽样; RA32 ratio: 2000/2000 有限重抽样; RA32source a: 2000/2000 有限重抽样; RA32source b: 2000/2000 有限重抽样; RA32source ratio: 2000/2000 有限重抽样.
\end{minipage}
\end{table}

\begin{table}[htbp]
\centering
\scriptsize
\caption{逐折模型验证}\label{tab:cv}
\setlength{\tabcolsep}{2pt}
\begin{tabularx}{\linewidth}{@{}>{\raggedright\arraybackslash}p{0.2400\linewidth}>{\raggedright\arraybackslash}p{0.1167\linewidth}>{\raggedright\arraybackslash}p{0.1167\linewidth}>{\raggedright\arraybackslash}p{0.1167\linewidth}>{\raggedright\arraybackslash}p{0.1167\linewidth}>{\raggedright\arraybackslash}p{0.1167\linewidth}>{\raggedright\arraybackslash}p{0.1167\linewidth}@{}}
\toprule
统计量/划分 & 折 & 有效/总数或n & 仿射 & 幂律 & 差值 & 胜出/有效 \\
\midrule
RA / 州 & 汇总 & 51/51 & -- & -- & -0.002018 & 26 \\
RA / 季 & 汇总 & 4/4 & -- & -- & 0.004626 & 3 \\
RA & 2022-23 & 365 & -2.378 & -2.369 & 0.008954 & 是 \\
RA & 2023-24 & 499 & -2.613 & -2.612 & 0.0008806 & 是 \\
RA & 2024-25 & 356 & -2.46 & -2.481 & -0.02077 & 是 \\
RA & 2025-26 & 528 & -2.346 & -2.324 & 0.0223 & 是 \\
G / 州 & 汇总 & 51/51 & -- & -- & -0.007369 & 22 \\
G / 季 & 汇总 & 4/4 & -- & -- & -0.004414 & 2 \\
G & 2022-23 & 367 & -0.5979 & -0.5942 & 0.003756 & 是 \\
G & 2023-24 & 503 & -1.201 & -1.192 & 0.008487 & 是 \\
G & 2024-25 & 357 & -0.7333 & -0.7584 & -0.02505 & 是 \\
G & 2025-26 & 529 & -0.7557 & -0.7641 & -0.008421 & 是 \\
RA32 / 州 & 汇总 & 51/51 & -- & -- & -0.002242 & 25 \\
RA32 / 季 & 汇总 & 4/4 & -- & -- & 0.004343 & 3 \\
RA32 & 2022-23 & 365 & -2.612 & -2.604 & 0.00805 & 是 \\
RA32 & 2023-24 & 499 & -2.83 & -2.83 & 0.0001619 & 是 \\
RA32 & 2024-25 & 356 & -2.691 & -2.712 & -0.02118 & 是 \\
RA32 & 2025-26 & 528 & -2.572 & -2.549 & 0.02294 & 是 \\
RA32source / 州 & 汇总 & 51/51 & -- & -- & -0.002824 & 25 \\
RA32source / 季 & 汇总 & 4/4 & -- & -- & 0.003577 & 2 \\
RA32source & 2022-23 & 365 & -2.701 & -2.695 & 0.005641 & 是 \\
RA32source & 2023-24 & 499 & -2.886 & -2.888 & -0.001735 & 是 \\
RA32source & 2024-25 & 356 & -2.77 & -2.792 & -0.02217 & 是 \\
RA32source & 2025-26 & 528 & -2.644 & -2.619 & 0.02453 & 是 \\
\bottomrule
\end{tabularx}
\par\smallskip
\begin{minipage}{\linewidth}\footnotesize\raggedright
汇总行列出有效/总折数及末列偏向仿射模型的折数；季节行列出测试样本量与收敛状态。评分为平均 L/mu + log(mu)，越低越好。差值=幂律减仿射，正值偏向仿射。各折重新拟合训练中心。训练集重叠，不报告独立折标准误或四季显著性检验。
\end{minipage}
\end{table}

\begin{table}[htbp]
\centering
\scriptsize
\caption{共享计数稀释估计、同时区间与模拟诊断}\label{tab:thinning}
\setlength{\tabcolsep}{2pt}
\begin{tabularx}{\linewidth}{@{}>{\raggedright\arraybackslash}p{0.2400\linewidth}>{\raggedright\arraybackslash}p{0.1400\linewidth}>{\raggedright\arraybackslash}p{0.1400\linewidth}>{\raggedright\arraybackslash}p{0.1400\linewidth}>{\raggedright\arraybackslash}p{0.1400\linewidth}>{\raggedright\arraybackslash}p{0.1400\linewidth}@{}}
\toprule
样本 & q/拟合 & 指标 & 估计值 [95\% CI] & 同时95\%区间 & 有限/B \\
\midrule
full & q=0.8 & b & 0.6747 [0.4905, 0.9034] & -- & 1000/1000 \\
full & q=0.8 & c & 0.4588 [0.3731, 0.5557] & [0.3381, 0.5795] & 1000/1000 \\
full & q=0.8 & delta\_\allowbreak{}a & -0.0003183 [-0.000786, 0.00005128] & [-0.0008833, 0.0002467] & 1000/1000 \\
full & q=0.8 & c\_\allowbreak{}over\_\allowbreak{}bP & 1.03 [0.8476, 1.242] & -- & 1000/1000 \\
full & q=0.6 & b & 0.8763 [0.6941, 1.103] & -- & 1000/1000 \\
full & q=0.6 & c & 0.4745 [0.4064, 0.5532] & [0.3751, 0.574] & 1000/1000 \\
full & q=0.6 & delta\_\allowbreak{}a & -0.000836 [-0.001963, -0.000006555] & [-0.00213, 0.0004583] & 1000/1000 \\
full & q=0.6 & c\_\allowbreak{}over\_\allowbreak{}bP & 1.065 [0.915, 1.241] & -- & 1000/1000 \\
full & q=0.5 & b & 1.057 [0.858, 1.275] & -- & 1000/1000 \\
full & q=0.5 & c & 0.4969 [0.4243, 0.5606] & [0.4047, 0.589] & 1000/1000 \\
full & q=0.5 & delta\_\allowbreak{}a & -0.001318 [-0.002744, -0.00009326] & [-0.003126, 0.0004903] & 1000/1000 \\
full & q=0.5 & c\_\allowbreak{}over\_\allowbreak{}bP & 1.115 [0.9621, 1.252] & -- & 1000/1000 \\
full & q=0.4 & b & 1.31 [1.112, 1.534] & -- & 1000/1000 \\
full & q=0.4 & c & 0.4997 [0.4397, 0.5611] & [0.4158, 0.5837] & 1000/1000 \\
full & q=0.4 & delta\_\allowbreak{}a & -0.001794 [-0.00376, -0.0001502] & [-0.004272, 0.0006837] & 1000/1000 \\
full & q=0.4 & c\_\allowbreak{}over\_\allowbreak{}bP & 1.122 [0.9903, 1.257] & -- & 1000/1000 \\
full & q=0.3 & b & 1.742 [1.503, 1.958] & -- & 1000/1000 \\
full & q=0.3 & c & 0.5067 [0.4482, 0.5539] & [0.4326, 0.5808] & 1000/1000 \\
full & q=0.3 & delta\_\allowbreak{}a & -0.002454 [-0.004852, -0.000159] & [-0.005728, 0.0008198] & 1000/1000 \\
full & q=0.3 & c\_\allowbreak{}over\_\allowbreak{}bP & 1.137 [1.008, 1.249] & -- & 1000/1000 \\
full & q=0.2 & b & 2.46 [2.221, 2.711] & -- & 1000/1000 \\
full & q=0.2 & c & 0.4749 [0.4292, 0.5229] & [0.4126, 0.5373] & 1000/1000 \\
full & q=0.2 & delta\_\allowbreak{}a & -0.002484 [-0.005751, 0.0005387] & [-0.006867, 0.001899] & 1000/1000 \\
full & q=0.2 & c\_\allowbreak{}over\_\allowbreak{}bP & 1.066 [0.9728, 1.173] & -- & 1000/1000 \\
full & q=0.1 & b & 3.983 [3.509, 4.379] & -- & 1000/1000 \\
full & q=0.1 & c & 0.3804 [0.3337, 0.4204] & [0.3234, 0.4373] & 1000/1000 \\
full & q=0.1 & delta\_\allowbreak{}a & 0.0007312 [-0.00305, 0.005984] & [-0.005595, 0.007057] & 1000/1000 \\
full & q=0.1 & c\_\allowbreak{}over\_\allowbreak{}bP & 0.8537 [0.7415, 0.9566] & -- & 1000/1000 \\
full & full & c & 0.4051 [0.3656, 0.4413] & -- & 1000/1000 \\
full & full & c\_\allowbreak{}over\_\allowbreak{}bP & 0.9093 [0.8101, 0.9995] & -- & 1000/1000 \\
full & full & delta\_\allowbreak{}a\_\allowbreak{}rate & -0.0001293 [-0.0006192, 0.0004699] & -- & 1000/1000 \\
full & noq01 & c & 0.4848 [0.4387, 0.5279] & -- & 1000/1000 \\
full & noq01 & c\_\allowbreak{}over\_\allowbreak{}bP & 1.088 [0.9936, 1.186] & -- & 1000/1000 \\
full & noq01 & delta\_\allowbreak{}a\_\allowbreak{}rate & -0.0008058 [-0.001703, 0.00005179] & -- & 1000/1000 \\
highCt50 & q=0.8 & b & 0.7797 [-0.007567, 1.573] & -- & 1000/1000 \\
highCt50 & q=0.8 & c & 0.442 [0.2771, 0.6691] & [0.1838, 0.7002] & 1000/1000 \\
highCt50 & q=0.8 & delta\_\allowbreak{}a & -0.00007119 [-0.0004223, 0.0001553] & [-0.0004703, 0.0003279] & 1000/1000 \\
highCt50 & q=0.8 & c\_\allowbreak{}over\_\allowbreak{}bP & 1.051 [0.6634, 1.606] & -- & 1000/1000 \\
highCt50 & q=0.6 & b & 0.969 [0.1971, 1.773] & -- & 1000/1000 \\
highCt50 & q=0.6 & c & 0.4497 [0.3339, 0.6318] & [0.256, 0.6433] & 1000/1000 \\
highCt50 & q=0.6 & delta\_\allowbreak{}a & -0.0001877 [-0.0009664, 0.0002652] & [-0.0009926, 0.0006171] & 1000/1000 \\
highCt50 & q=0.6 & c\_\allowbreak{}over\_\allowbreak{}bP & 1.069 [0.8015, 1.51] & -- & 1000/1000 \\
highCt50 & q=0.5 & b & 1.126 [0.3593, 1.891] & -- & 1000/1000 \\
highCt50 & q=0.5 & c & 0.4567 [0.3405, 0.6213] & [0.273, 0.6404] & 1000/1000 \\
highCt50 & q=0.5 & delta\_\allowbreak{}a & -0.0002913 [-0.001459, 0.0003529] & [-0.001456, 0.000873] & 1000/1000 \\
highCt50 & q=0.5 & c\_\allowbreak{}over\_\allowbreak{}bP & 1.086 [0.8187, 1.476] & -- & 1000/1000 \\
highCt50 & q=0.4 & b & 1.351 [0.6322, 2.134] & -- & 1000/1000 \\
highCt50 & q=0.4 & c & 0.4548 [0.3568, 0.593] & [0.2972, 0.6124] & 1000/1000 \\
highCt50 & q=0.4 & delta\_\allowbreak{}a & -0.0004078 [-0.001854, 0.0004131] & [-0.001938, 0.001122] & 1000/1000 \\
highCt50 & q=0.4 & c\_\allowbreak{}over\_\allowbreak{}bP & 1.081 [0.8483, 1.413] & -- & 1000/1000 \\
highCt50 & q=0.3 & b & 1.723 [1.005, 2.495] & -- & 1000/1000 \\
highCt50 & q=0.3 & c & 0.4515 [0.3668, 0.5711] & [0.3176, 0.5854] & 1000/1000 \\
highCt50 & q=0.3 & delta\_\allowbreak{}a & -0.0005649 [-0.00247, 0.0006027] & [-0.002578, 0.001448] & 1000/1000 \\
highCt50 & q=0.3 & c\_\allowbreak{}over\_\allowbreak{}bP & 1.073 [0.8687, 1.355] & -- & 1000/1000 \\
highCt50 & q=0.2 & b & 2.522 [1.77, 3.276] & -- & 1000/1000 \\
highCt50 & q=0.2 & c & 0.4633 [0.3813, 0.5611] & [0.3435, 0.5831] & 1000/1000 \\
highCt50 & q=0.2 & delta\_\allowbreak{}a & -0.0009427 [-0.00351, 0.0008025] & [-0.003822, 0.001937] & 1000/1000 \\
highCt50 & q=0.2 & c\_\allowbreak{}over\_\allowbreak{}bP & 1.102 [0.9074, 1.325] & -- & 1000/1000 \\
highCt50 & q=0.1 & b & 4.97 [4.216, 5.814] & -- & 999/1000 \\
highCt50 & q=0.1 & c & 0.4778 [0.4127, 0.5552] & [0.3835, 0.5721] & 999/1000 \\
highCt50 & q=0.1 & delta\_\allowbreak{}a & -0.001594 [-0.005615, 0.001213] & [-0.006149, 0.002962] & 999/1000 \\
highCt50 & q=0.1 & c\_\allowbreak{}over\_\allowbreak{}bP & 1.136 [0.9843, 1.323] & -- & 999/1000 \\
highCt50 & full & c & 0.4734 [0.409, 0.5503] & -- & 999/1000 \\
highCt50 & full & c\_\allowbreak{}over\_\allowbreak{}bP & 1.126 [0.9702, 1.316] & -- & 999/1000 \\
highCt50 & full & delta\_\allowbreak{}a\_\allowbreak{}rate & -0.0001928 [-0.0006891, 0.0001441] & -- & 999/1000 \\
highCt50 & noq01 & c & 0.4594 [0.3808, 0.5603] & -- & 1000/1000 \\
highCt50 & noq01 & c\_\allowbreak{}over\_\allowbreak{}bP & 1.092 [0.9026, 1.327] & -- & 1000/1000 \\
highCt50 & noq01 & delta\_\allowbreak{}a\_\allowbreak{}rate & -0.0002434 [-0.0009563, 0.0002083] & -- & 1000/1000 \\
both50 & q=0.8 & b & 0.5473 [-0.2021, 1.391] & -- & 1000/1000 \\
both50 & q=0.8 & c & 0.4994 [0.3073, 0.7827] & [0.1979, 0.8009] & 1000/1000 \\
both50 & q=0.8 & delta\_\allowbreak{}a & -0.0001674 [-0.0006684, 0.0001096] & [-0.0006864, 0.0003517] & 1000/1000 \\
both50 & q=0.8 & c\_\allowbreak{}over\_\allowbreak{}bP & 1.194 [0.722, 1.896] & -- & 1000/1000 \\
both50 & q=0.6 & b & 0.754 [0.02843, 1.586] & -- & 1000/1000 \\
both50 & q=0.6 & c & 0.4973 [0.3582, 0.7249] & [0.2629, 0.7317] & 1000/1000 \\
both50 & q=0.6 & delta\_\allowbreak{}a & -0.0004007 [-0.001555, 0.000186] & [-0.001503, 0.0007021] & 1000/1000 \\
both50 & q=0.6 & c\_\allowbreak{}over\_\allowbreak{}bP & 1.189 [0.854, 1.751] & -- & 1000/1000 \\
both50 & q=0.5 & b & 0.9288 [0.1947, 1.748] & -- & 1000/1000 \\
both50 & q=0.5 & c & 0.5063 [0.3635, 0.7058] & [0.2897, 0.723] & 1000/1000 \\
both50 & q=0.5 & delta\_\allowbreak{}a & -0.0006083 [-0.00219, 0.0001859] & [-0.002149, 0.0009326] & 1000/1000 \\
both50 & q=0.5 & c\_\allowbreak{}over\_\allowbreak{}bP & 1.21 [0.8863, 1.69] & -- & 1000/1000 \\
both50 & q=0.4 & b & 1.161 [0.4133, 1.977] & -- & 1000/1000 \\
both50 & q=0.4 & c & 0.4921 [0.3785, 0.6457] & [0.3095, 0.6746] & 1000/1000 \\
both50 & q=0.4 & delta\_\allowbreak{}a & -0.00078 [-0.002694, 0.0002525] & [-0.002738, 0.001178] & 1000/1000 \\
both50 & q=0.4 & c\_\allowbreak{}over\_\allowbreak{}bP & 1.176 [0.9064, 1.548] & -- & 1000/1000 \\
both50 & q=0.3 & b & 1.554 [0.8611, 2.363] & -- & 1000/1000 \\
both50 & q=0.3 & c & 0.4849 [0.3882, 0.6214] & [0.3325, 0.6373] & 1000/1000 \\
both50 & q=0.3 & delta\_\allowbreak{}a & -0.001058 [-0.003465, 0.0003471] & [-0.003573, 0.001457] & 1000/1000 \\
both50 & q=0.3 & c\_\allowbreak{}over\_\allowbreak{}bP & 1.159 [0.9271, 1.485] & -- & 1000/1000 \\
both50 & q=0.2 & b & 2.382 [1.63, 3.167] & -- & 1000/1000 \\
both50 & q=0.2 & c & 0.4899 [0.3955, 0.5995] & [0.3589, 0.6209] & 1000/1000 \\
both50 & q=0.2 & delta\_\allowbreak{}a & -0.001588 [-0.004831, 0.0005575] & [-0.005049, 0.001874] & 1000/1000 \\
both50 & q=0.2 & c\_\allowbreak{}over\_\allowbreak{}bP & 1.171 [0.9456, 1.432] & -- & 1000/1000 \\
both50 & q=0.1 & b & 4.863 [4.1, 5.734] & -- & 999/1000 \\
both50 & q=0.1 & c & 0.4934 [0.4221, 0.5729] & [0.3956, 0.5913] & 999/1000 \\
both50 & q=0.1 & delta\_\allowbreak{}a & -0.002397 [-0.007015, 0.000824] & [-0.007548, 0.002755] & 999/1000 \\
both50 & q=0.1 & c\_\allowbreak{}over\_\allowbreak{}bP & 1.18 [1.005, 1.371] & -- & 999/1000 \\
both50 & full & c & 0.4926 [0.42, 0.5742] & -- & 999/1000 \\
both50 & full & c\_\allowbreak{}over\_\allowbreak{}bP & 1.178 [1.004, 1.379] & -- & 999/1000 \\
both50 & full & delta\_\allowbreak{}a\_\allowbreak{}rate & -0.0003058 [-0.000896, 0.00009878] & -- & 999/1000 \\
both50 & noq01 & c & 0.4898 [0.3971, 0.6072] & -- & 1000/1000 \\
both50 & noq01 & c\_\allowbreak{}over\_\allowbreak{}bP & 1.171 [0.9533, 1.462] & -- & 1000/1000 \\
both50 & noq01 & delta\_\allowbreak{}a\_\allowbreak{}rate & -0.0004328 [-0.001331, 0.0001204] & -- & 1000/1000 \\
\bottomrule
\end{tabularx}
\par\smallskip
\begin{minipage}{\linewidth}\footnotesize\raggedright
逐q区间为点态区间；单列包络为各子集中跨q的同时区间。full使用全部q；noq01排除q=0.1。highCt50要求原当期计数>=50；both50要求原相邻两期计数均>=50。c为每(1/q-1)的斜率增量；bP为代入式泊松参照。 B=1000; R(point)=400; R(inner, actual)=50; precision criterion: MC SE at most 0.1 times the point estimate, applied to the five monitored response summaries; max finite MC SE/|point| over all stored metrics=68.44; max finite MC SE/sampling SE over all stored metrics=1.74; undefined MC ratios=0/219; MC status=tolerance\_\allowbreak{}met. Point fits 8400/8400; bootstrap fits 1049998/1050000; failed bootstrap averages=2. Any failed inner fit makes that q average undefined; no silent averaging over successful fits. full: n=1748; bP=0.4455 [0.4282, 0.4658]. highCt50: n=828; bP=0.4206 [0.4032, 0.4387]. both50: n=821; bP=0.4183 [0.4008, 0.4359].
\end{minipage}
\end{table}

\begin{table}[htbp]
\centering
\scriptsize
\caption{R\_\allowbreak{}A 样本（n=1748）的固定效应：配对比值、周块敏感性与收敛}\label{tab:fe}
\setlength{\tabcolsep}{2pt}
\begin{tabularx}{\linewidth}{@{}>{\raggedright\arraybackslash}p{0.2400\linewidth}>{\raggedright\arraybackslash}p{0.1750\linewidth}>{\raggedright\arraybackslash}p{0.1750\linewidth}>{\raggedright\arraybackslash}p{0.1750\linewidth}>{\raggedright\arraybackslash}p{0.1750\linewidth}@{}}
\toprule
周块 & 参数 & 估计值 [95\% CI] & 有限/B & 非正比例/诊断 \\
\midrule
L=1 & Pooled a & 0.01495 [0.009976, 0.02114] & 1000/1000 & 0 \\
L=1 & Pooled b & 0.56 [0.3682, 0.7707] & 1000/1000 & 0 \\
L=1 & Pooled b/a & 37.44 [18.64, 73.92] & 1000/1000 & 0 \\
L=1 & State FE b & 0.3794 [0.1082, 0.6403] & 999/1000 & 0 \\
L=1 & State FE mean a & 0.02192 [0.01289, 0.03502] & 999/1000 & 0 \\
L=1 & State FE b/mean a & 17.3 [3.315, 47.8] & 999/1000 & 0 \\
L=1 & State FE median a & 0.01751 [0.01037, 0.03014] & 999/1000 & 0 \\
L=1 & State FE b/median a & 21.67 [3.977, 57.85] & 999/1000 & 0 \\
L=1 & State FE row-weighted a & 0.02061 [0.01339, 0.03114] & 999/1000 & 0 \\
L=1 & State FE row-weighted ratio & 18.41 [3.705, 45.56] & 999/1000 & 0 \\
L=1 & State FE count-weighted a & 0.01542 [0.01056, 0.02288] & 999/1000 & 0 \\
L=1 & State FE count-weighted ratio & 24.61 [5.265, 55.65] & 999/1000 & 0 \\
L=1 & True D\_\allowbreak{}t b & 0.3691 [0.1153, 0.6145] & 999/1000 & 0 \\
L=1 & True D\_\allowbreak{}t mean a & 0.02084 [0.01218, 0.03429] & 999/1000 & 0 \\
L=1 & True D\_\allowbreak{}t ratio & 17.71 [3.267, 47.81] & 999/1000 & 0 \\
L=1 & Paired pooled/FE ratio & 2.164 [1.04, 7.745] & 999/1000 & 0 \\
L=2 & Pooled a & 0.01495 [0.009493, 0.02077] & 1000/1000 & 0 \\
L=2 & Pooled b & 0.56 [0.371, 0.763] & 1000/1000 & 0 \\
L=2 & Pooled b/a & 37.44 [18.73, 74.77] & 1000/1000 & 0 \\
L=2 & State FE b & 0.3794 [0.109, 0.6611] & 1000/1000 & 0 \\
L=2 & State FE mean a & 0.02192 [0.01288, 0.03581] & 1000/1000 & 0 \\
L=2 & State FE b/mean a & 17.3 [3.453, 50.16] & 1000/1000 & 0 \\
L=2 & State FE median a & 0.01751 [0.01024, 0.03013] & 1000/1000 & 0 \\
L=2 & State FE b/median a & 21.67 [3.984, 60.15] & 1000/1000 & 0 \\
L=2 & State FE row-weighted a & 0.02061 [0.013, 0.03107] & 1000/1000 & 0 \\
L=2 & State FE row-weighted ratio & 18.41 [3.819, 48.17] & 1000/1000 & 0 \\
L=2 & State FE count-weighted a & 0.01542 [0.0102, 0.02293] & 1000/1000 & 0 \\
L=2 & State FE count-weighted ratio & 24.61 [5.111, 57.08] & 1000/1000 & 0 \\
L=2 & True D\_\allowbreak{}t b & 0.3691 [0.1082, 0.612] & 1000/1000 & 0.001 \\
L=2 & True D\_\allowbreak{}t mean a & 0.02084 [0.01168, 0.03555] & 1000/1000 & 0 \\
L=2 & True D\_\allowbreak{}t ratio & 17.71 [3.26, 51.13] & 1000/1000 & 0.001 \\
L=2 & Paired pooled/FE ratio & 2.164 [0.964, 8.483] & 1000/1000 & 0 \\
L=3 & Pooled a & 0.01495 [0.01001, 0.02089] & 1000/1000 & 0 \\
L=3 & Pooled b & 0.56 [0.3851, 0.7603] & 1000/1000 & 0 \\
L=3 & Pooled b/a & 37.44 [19.39, 68.59] & 1000/1000 & 0 \\
L=3 & State FE b & 0.3794 [0.111, 0.6669] & 1000/1000 & 0.001 \\
L=3 & State FE mean a & 0.02192 [0.01247, 0.03692] & 1000/1000 & 0 \\
L=3 & State FE b/mean a & 17.3 [3.051, 50.11] & 1000/1000 & 0.001 \\
L=3 & State FE median a & 0.01751 [0.01068, 0.03045] & 1000/1000 & 0 \\
L=3 & State FE b/median a & 21.67 [3.982, 58.24] & 1000/1000 & 0.001 \\
L=3 & State FE row-weighted a & 0.02061 [0.01293, 0.0317] & 1000/1000 & 0 \\
L=3 & State FE row-weighted ratio & 18.41 [3.407, 47.65] & 1000/1000 & 0.001 \\
L=3 & State FE count-weighted a & 0.01542 [0.01076, 0.02338] & 1000/1000 & 0 \\
L=3 & State FE count-weighted ratio & 24.61 [4.841, 57.62] & 1000/1000 & 0.001 \\
L=3 & True D\_\allowbreak{}t b & 0.3691 [0.1072, 0.6252] & 1000/1000 & 0.001 \\
L=3 & True D\_\allowbreak{}t mean a & 0.02084 [0.0116, 0.03629] & 1000/1000 & 0 \\
L=3 & True D\_\allowbreak{}t ratio & 17.71 [3.068, 51.41] & 1000/1000 & 0.001 \\
L=3 & Paired pooled/FE ratio & 2.164 [0.9866, 8.87] & 999/1000 & 0 \\
点估计诊断 & week\_\allowbreak{}only & b=0.5258 & 无区间 & group\_\allowbreak{}gamma\_\allowbreak{}objective\_\allowbreak{}fallback \\
点估计诊断 & state\_\allowbreak{}and\_\allowbreak{}week & 不可用 & 无区间 & iteration\_\allowbreak{}limit \\
\bottomrule
\end{tabularx}
\par\smallskip
\begin{minipage}{\linewidth}\footnotesize\raggedright
样本为 \$R\_\allowbreak{}A\$ 规格的 1748 个州周、51 个辖区；本表不使用 \$G\$ 样本。区间保留有符号重抽样；有限/B<1时，区间以有限估计为条件。末列为全部请求重抽样中的非正比例。FE以抽中的州出现次数为目标；合并/FE比较在重抽样内配对。L为季内周块长度。仅周和州加周拟合为点估计诊断，失败时无估计或区间。本表合并行使用本表自身的 1000 次自助抽样，可能与表 S1 的 2000 次主区间略有差异。 L=1 pooled: 1000 converged, 0 failed; \{\}. L=1 state: 999 converged, 1 failed; \{"gamma\_\allowbreak{}objective\_\allowbreak{}fallback\_\allowbreak{}failed": 1\}. L=1 state\_\allowbreak{}D: 999 converged, 1 failed; \{"gamma\_\allowbreak{}objective\_\allowbreak{}fallback\_\allowbreak{}failed": 1\}. L=2 pooled: 1000 converged, 0 failed; \{\}. L=2 state: 1000 converged, 0 failed; \{\}. L=2 state\_\allowbreak{}D: 1000 converged, 0 failed; \{\}. L=3 pooled: 1000 converged, 0 failed; \{\}. L=3 state: 1000 converged, 0 failed; \{\}. L=3 state\_\allowbreak{}D: 1000 converged, 0 failed; \{\}.
\end{minipage}
\end{table}

\begin{table}[htbp]
\centering
\scriptsize
\caption{采用共同损失分母的中心化阶段}\label{tab:centering}
\setlength{\tabcolsep}{2pt}
\begin{tabularx}{\linewidth}{@{}>{\raggedright\arraybackslash}p{0.2400\linewidth}>{\raggedright\arraybackslash}p{0.1750\linewidth}>{\raggedright\arraybackslash}p{0.1750\linewidth}>{\raggedright\arraybackslash}p{0.1750\linewidth}>{\raggedright\arraybackslash}p{0.1750\linewidth}@{}}
\toprule
统计量/中心 & 合并均值 & a / 变化 [95\% CI] & b [95\% CI] & 有限/B \\
\midrule
RA / pooled & 1.206 & 0.01495 [0.01016, 0.02083] & 0.56 [0.3625, 0.7639] & 1000/1000 \\
RA / season & 1.206 & 0.01471 [0.009507, 0.01966] & 0.5391 [0.3602, 0.7205] & 1000/1000 \\
RA / season\_\allowbreak{}linear & 1.206 & 0.0154 [0.008972, 0.01901] & 0.3936 [0.2717, 0.5763] & 1000/1000 \\
RA / week & 1.206 & 0.009766 [0.007103, 0.01357] & 0.4581 [0.3124, 0.5916] & 1000/1000 \\
RA / week\_\allowbreak{}leave\_\allowbreak{}state\_\allowbreak{}out & 1.206 & 0.01016 [0.007654, 0.01481] & 0.4841 [0.3492, 0.6636] & 1000/1000 \\
RA / week\_\allowbreak{}change & 1.206 & -0.005189 [-0.01015, -0.0006938] & -- & 1000/1000 \\
RA / week\_\allowbreak{}reduction\_\allowbreak{}fraction & 1.206 & 34.7 [5.967, 51.9] & -- & 1000/1000 \\
G / pooled & 1.724 & 0.08401 [0.05163, 0.1222] & 1.971 [0.9987, 2.986] & 1000/1000 \\
G / season & 1.724 & 0.08163 [0.0504, 0.1156] & 1.842 [0.9591, 2.602] & 1000/1000 \\
G / season\_\allowbreak{}linear & 1.724 & 0.07864 [0.04808, 0.09897] & 1.249 [0.6423, 1.988] & 1000/1000 \\
G / week & 1.724 & 0.05607 [0.03962, 0.07812] & 1.401 [0.7134, 1.912] & 1000/1000 \\
G / week\_\allowbreak{}leave\_\allowbreak{}state\_\allowbreak{}out & 1.724 & 0.05843 [0.04308, 0.08505] & 1.477 [0.8089, 2.138] & 1000/1000 \\
G / week\_\allowbreak{}change & 1.724 & -0.02794 [-0.05687, -0.003708] & -- & 1000/1000 \\
G / week\_\allowbreak{}reduction\_\allowbreak{}fraction & 1.724 & 33.25 [6.805, 51.63] & -- & 1000/1000 \\
\bottomrule
\end{tabularx}
\par\smallskip
\begin{minipage}{\linewidth}\footnotesize\raggedright
同一统计量的五种中心均使用相同合并均值作为相对损失分母；每次重抽样用权重重新拟合该分母。week\_\allowbreak{}leave\_\allowbreak{}state\_\allowbreak{}out的周中心排除本州。减少比例为100*(a\_\allowbreak{}pooled-a\_\allowbreak{}week)/a\_\allowbreak{}pooled，单位为百分比；这些是描述性变化，并非可加总方差分量。本表区间使用自身的 1000 次自助抽样，可能与表 S1 的 2000 次主区间略有差异。 RA/pooled: 1000/1000 有限重抽样; RA/season: 1000/1000 有限重抽样; RA/season\_\allowbreak{}linear: 1000/1000 有限重抽样; RA/week: 1000/1000 有限重抽样; RA/week\_\allowbreak{}leave\_\allowbreak{}state\_\allowbreak{}out: 1000/1000 有限重抽样; G/pooled: 1000/1000 有限重抽样; G/season: 1000/1000 有限重抽样; G/season\_\allowbreak{}linear: 1000/1000 有限重抽样; G/week: 1000/1000 有限重抽样; G/week\_\allowbreak{}leave\_\allowbreak{}state\_\allowbreak{}out: 1000/1000 有限重抽样.
\end{minipage}
\end{table}

\begin{table}[htbp]
\centering
\scriptsize
\caption{全部声明情景的条件校准}\label{tab:calibration}
\setlength{\tabcolsep}{2pt}
\begin{tabularx}{\linewidth}{@{}>{\raggedright\arraybackslash}p{0.2400\linewidth}>{\raggedright\arraybackslash}p{0.1400\linewidth}>{\raggedright\arraybackslash}p{0.1400\linewidth}>{\raggedright\arraybackslash}p{0.1400\linewidth}>{\raggedright\arraybackslash}p{0.1400\linewidth}>{\raggedright\arraybackslash}p{0.1400\linewidth}@{}}
\toprule
情景/统计量 & v [95\% CI] & 有限/B & 非正/非有限 & 超网格/斜率<=0 & 失败拟合/全部 \\
\midrule
delay\_\allowbreak{}long / RA & 0.04587 [0.02996, 0.06461] & 2000/2000 & 0/0 & 0/0 & 0/1000 \\
delay\_\allowbreak{}long / G & 0.04686 [0.02666, 0.06952] & 2000/2000 & 0/0 & 3/0 & 0/1000 \\
delay\_\allowbreak{}short / RA & 0.03863 [0.02531, 0.05527] & 2000/2000 & 0/0 & 0/0 & 0/1000 \\
delay\_\allowbreak{}short / G & 0.04161 [0.02382, 0.06314] & 2000/2000 & 0/0 & 0/0 & 0/1000 \\
joint\_\allowbreak{}gi3.2\_\allowbreak{}pn0.25\_\allowbreak{}obs1 / RA & 0.02882 [0.01937, 0.04107] & 2000/2000 & 0/0 & 0/0 & 0/1000 \\
joint\_\allowbreak{}gi3.2\_\allowbreak{}pn0.25\_\allowbreak{}obs1 / G & 0.02967 [0.01709, 0.04456] & 2000/2000 & 0/0 & 0/0 & 0/1000 \\
joint\_\allowbreak{}gi3.2\_\allowbreak{}pn0.25\_\allowbreak{}obs1.25 / RA & 0.02844 [0.01853, 0.03995] & 2000/2000 & 0/0 & 0/0 & 0/1000 \\
joint\_\allowbreak{}gi3.2\_\allowbreak{}pn0.25\_\allowbreak{}obs1.25 / G & 0.0294 [0.01681, 0.04331] & 2000/2000 & 0/0 & 0/0 & 0/1000 \\
joint\_\allowbreak{}gi3.2\_\allowbreak{}pn0\_\allowbreak{}obs1 / RA & 0.03178 [0.02085, 0.04488] & 2000/2000 & 0/0 & 0/0 & 0/1000 \\
joint\_\allowbreak{}gi3.2\_\allowbreak{}pn0\_\allowbreak{}obs1 / G & 0.03369 [0.01924, 0.05045] & 2000/2000 & 0/0 & 0/0 & 0/1000 \\
joint\_\allowbreak{}gi3.2\_\allowbreak{}pn0\_\allowbreak{}obs1.25 / RA & 0.03128 [0.02019, 0.04538] & 2000/2000 & 0/0 & 0/0 & 0/1000 \\
joint\_\allowbreak{}gi3.2\_\allowbreak{}pn0\_\allowbreak{}obs1.25 / G & 0.03348 [0.01895, 0.05184] & 2000/2000 & 0/0 & 0/0 & 0/1000 \\
joint\_\allowbreak{}gi3.6\_\allowbreak{}pn0.25\_\allowbreak{}obs1 / RA & 0.03764 [0.02444, 0.05359] & 2000/2000 & 0/0 & 0/0 & 0/1000 \\
joint\_\allowbreak{}gi3.6\_\allowbreak{}pn0.25\_\allowbreak{}obs1 / G & 0.03902 [0.02247, 0.05689] & 2000/2000 & 0/0 & 0/0 & 0/1000 \\
joint\_\allowbreak{}gi3.6\_\allowbreak{}pn0.25\_\allowbreak{}obs1.25 / RA & 0.03638 [0.02354, 0.0519] & 2000/2000 & 0/0 & 0/0 & 0/1000 \\
joint\_\allowbreak{}gi3.6\_\allowbreak{}pn0.25\_\allowbreak{}obs1.25 / G & 0.03785 [0.0216, 0.05704] & 2000/2000 & 0/0 & 0/0 & 0/1000 \\
joint\_\allowbreak{}gi3.6\_\allowbreak{}pn0\_\allowbreak{}obs1 / RA & 0.04002 [0.02558, 0.05767] & 2000/2000 & 0/0 & 0/0 & 0/1000 \\
joint\_\allowbreak{}gi3.6\_\allowbreak{}pn0\_\allowbreak{}obs1 / G & 0.04282 [0.02407, 0.06368] & 2000/2000 & 0/0 & 0/0 & 0/1000 \\
joint\_\allowbreak{}gi3.6\_\allowbreak{}pn0\_\allowbreak{}obs1.25 / RA & 0.04035 [0.02648, 0.0569] & 2000/2000 & 0/0 & 0/0 & 0/1000 \\
joint\_\allowbreak{}gi3.6\_\allowbreak{}pn0\_\allowbreak{}obs1.25 / G & 0.04356 [0.02576, 0.064] & 2000/2000 & 0/0 & 0/0 & 0/1000 \\
k01 / RA & 0.04023 [0.02553, 0.05736] & 2000/2000 & 0/0 & 0/0 & 0/1000 \\
k01 / G & 0.04289 [0.02475, 0.0645] & 2000/2000 & 0/0 & 0/0 & 0/1000 \\
matched32\_\allowbreak{}binomial / RA & 0.03108 [0.02037, 0.04399] & 2000/2000 & 0/0 & 0/0 & 0/1000 \\
matched32\_\allowbreak{}binomial / G & 0.03287 [0.0186, 0.0489] & 2000/2000 & 0/0 & 0/0 & 0/1000 \\
matched32\_\allowbreak{}source\_\allowbreak{}sd21\_\allowbreak{}binomial / RA & 0.02884 [0.01908, 0.04027] & 2000/2000 & 0/0 & 0/0 & 0/1000 \\
matched32\_\allowbreak{}source\_\allowbreak{}sd21\_\allowbreak{}binomial / G & 0.031 [0.01794, 0.04533] & 2000/2000 & 0/0 & 0/0 & 0/1000 \\
matched36\_\allowbreak{}binomial / RA & 0.04 [0.02623, 0.0562] & 2000/2000 & 0/0 & 0/0 & 0/1000 \\
matched36\_\allowbreak{}binomial / G & 0.04298 [0.02466, 0.06284] & 2000/2000 & 0/0 & 0/0 & 0/1000 \\
national\_\allowbreak{}persistence25 / RA & 0.0369 [0.0238, 0.05109] & 2000/2000 & 0/0 & 0/0 & 0/1000 \\
national\_\allowbreak{}persistence25 / G & 0.03817 [0.02185, 0.05643] & 2000/2000 & 0/0 & 0/0 & 0/1000 \\
national\_\allowbreak{}share0 / RA & 0.03935 [0.0255, 0.05576] & 2000/2000 & 0/0 & 0/0 & 0/1000 \\
national\_\allowbreak{}share0 / G & 0.0416 [0.02393, 0.06162] & 2000/2000 & 0/0 & 0/0 & 0/1000 \\
national\_\allowbreak{}share65 / RA & 0.04014 [0.02576, 0.05685] & 2000/2000 & 0/0 & 0/0 & 0/1000 \\
national\_\allowbreak{}share65 / G & 0.04327 [0.02461, 0.06553] & 2000/2000 & 0/0 & 0/0 & 0/1000 \\
reporting\_\allowbreak{}shock / RA & 0.01954 [0.006213, 0.0358] & 2000/2000 & 0/0 & 0/0 & 0/1000 \\
reporting\_\allowbreak{}shock / G & 0.02903 [0.0117, 0.04918] & 2000/2000 & 0/0 & 0/0 & 0/1000 \\
rho\_\allowbreak{}h005 / RA & 0.03991 [0.02639, 0.05628] & 2000/2000 & 0/0 & 0/0 & 0/1000 \\
rho\_\allowbreak{}h005 / G & 0.04283 [0.02548, 0.06297] & 2000/2000 & 0/0 & 0/0 & 0/1000 \\
rho\_\allowbreak{}h02 / RA & 0.04029 [0.02632, 0.05694] & 2000/2000 & 0/0 & 0/0 & 0/1000 \\
rho\_\allowbreak{}h02 / G & 0.04317 [0.02494, 0.06466] & 2000/2000 & 0/0 & 0/0 & 0/1000 \\
state\_\allowbreak{}persistence50 / RA & 0.03244 [0.02095, 0.04559] & 2000/2000 & 0/0 & 0/0 & 0/1000 \\
state\_\allowbreak{}persistence50 / G & 0.03137 [0.01788, 0.0467] & 2000/2000 & 0/0 & 0/0 & 0/1000 \\
true2.2\_\allowbreak{}analyst36 / RA & 0.0146 [0.009442, 0.02045] & 2000/2000 & 0/0 & 0/0 & 0/1000 \\
true2.2\_\allowbreak{}analyst36 / G & 0.01473 [0.008358, 0.02141] & 2000/2000 & 0/0 & 0/0 & 0/1000 \\
true2.9\_\allowbreak{}analyst36 / RA & 0.02551 [0.0168, 0.03572] & 2000/2000 & 0/0 & 0/0 & 0/1000 \\
true2.9\_\allowbreak{}analyst36 / G & 0.02679 [0.01524, 0.0401] & 2000/2000 & 0/0 & 0/0 & 0/1000 \\
true4.3\_\allowbreak{}analyst36 / RA & 0.05807 [0.03763, 0.08382] & 2000/2000 & 0/0 & 27/0 & 0/1000 \\
true4.3\_\allowbreak{}analyst36 / G & 0.06312 [0.0366, 0.0955] & 2000/2000 & 0/0 & 78/0 & 0/1000 \\
\bottomrule
\end{tabularx}
\par\smallskip
\begin{minipage}{\linewidth}\footnotesize\raggedright
v为给定模型下、施加下界之前的周边际方差，不是已识别的生物学方差或上界。区间取全部有限有符号重抽样的2.5和97.5百分位；不作仅保留正值的筛选或截断。各列分别列出非正/非有限反演次数、超网格/非正斜率次数，以及失败/尝试的模拟拟合次数。 delay\_\allowbreak{}long: grid=[0.0, 0.02, 0.04, 0.06, 0.09]; RA point outside=False, nonpositive=False, max |line residual|/MC SE=0.8913; G point outside=False, nonpositive=False, max |line residual|/MC SE=0.5474. delay\_\allowbreak{}short: grid=[0.0, 0.02, 0.04, 0.06, 0.09]; RA point outside=False, nonpositive=False, max |line residual|/MC SE=1.781; G point outside=False, nonpositive=False, max |line residual|/MC SE=1.162. joint\_\allowbreak{}gi3.2\_\allowbreak{}pn0.25\_\allowbreak{}obs1: grid=[0.0, 0.02, 0.04, 0.06, 0.09]; RA point outside=False, nonpositive=False, max |line residual|/MC SE=0.974; G point outside=False, nonpositive=False, max |line residual|/MC SE=1.154. joint\_\allowbreak{}gi3.2\_\allowbreak{}pn0.25\_\allowbreak{}obs1.25: grid=[0.0, 0.02, 0.04, 0.06, 0.09]; RA point outside=False, nonpositive=False, max |line residual|/MC SE=1.557; G point outside=False, nonpositive=False, max |line residual|/MC SE=1.668. joint\_\allowbreak{}gi3.2\_\allowbreak{}pn0\_\allowbreak{}obs1: grid=[0.0, 0.02, 0.04, 0.06, 0.09]; RA point outside=False, nonpositive=False, max |line residual|/MC SE=1.13; G point outside=False, nonpositive=False, max |line residual|/MC SE=1.116. joint\_\allowbreak{}gi3.2\_\allowbreak{}pn0\_\allowbreak{}obs1.25: grid=[0.0, 0.02, 0.04, 0.06, 0.09]; RA point outside=False, nonpositive=False, max |line residual|/MC SE=1.369; G point outside=False, nonpositive=False, max |line residual|/MC SE=1.917. joint\_\allowbreak{}gi3.6\_\allowbreak{}pn0.25\_\allowbreak{}obs1: grid=[0.0, 0.02, 0.04, 0.06, 0.09]; RA point outside=False, nonpositive=False, max |line residual|/MC SE=1.666; G point outside=False, nonpositive=False, max |line residual|/MC SE=1.308. joint\_\allowbreak{}gi3.6\_\allowbreak{}pn0.25\_\allowbreak{}obs1.25: grid=[0.0, 0.02, 0.04, 0.06, 0.09]; RA point outside=False, nonpositive=False, max |line residual|/MC SE=1.255; G point outside=False, nonpositive=False, max |line residual|/MC SE=1.081. joint\_\allowbreak{}gi3.6\_\allowbreak{}pn0\_\allowbreak{}obs1: grid=[0.0, 0.02, 0.04, 0.06, 0.09]; RA point outside=False, nonpositive=False, max |line residual|/MC SE=1.596; G point outside=False, nonpositive=False, max |line residual|/MC SE=1.363. joint\_\allowbreak{}gi3.6\_\allowbreak{}pn0\_\allowbreak{}obs1.25: grid=[0.0, 0.02, 0.04, 0.06, 0.09]; RA point outside=False, nonpositive=False, max |line residual|/MC SE=0.9047; G point outside=False, nonpositive=False, max |line residual|/MC SE=0.9728. k01: grid=[0.0, 0.02, 0.04, 0.06, 0.09]; RA point outside=False, nonpositive=False, max |line residual|/MC SE=1.074; G point outside=False, nonpositive=False, max |line residual|/MC SE=1.429. matched32\_\allowbreak{}binomial: grid=[0.0, 0.02, 0.04, 0.06, 0.09]; RA point outside=False, nonpositive=False, max |line residual|/MC SE=0.4468; G point outside=False, nonpositive=False, max |line residual|/MC SE=0.4955. matched32\_\allowbreak{}source\_\allowbreak{}sd21\_\allowbreak{}binomial: grid=[0.0, 0.02, 0.04, 0.06, 0.09]; RA point outside=False, nonpositive=False, max |line residual|/MC SE=1.05; G point outside=False, nonpositive=False, max |line residual|/MC SE=0.9843. matched36\_\allowbreak{}binomial: grid=[0.0, 0.02, 0.04, 0.06, 0.09]; RA point outside=False, nonpositive=False, max |line residual|/MC SE=2.388; G point outside=False, nonpositive=False, max |line residual|/MC SE=1.795. national\_\allowbreak{}persistence25: grid=[0.0, 0.02, 0.04, 0.06, 0.09]; RA point outside=False, nonpositive=False, max |line residual|/MC SE=2.246; G point outside=False, nonpositive=False, max |line residual|/MC SE=2.37. national\_\allowbreak{}share0: grid=[0.0, 0.02, 0.04, 0.06, 0.09]; RA point outside=False, nonpositive=False, max |line residual|/MC SE=0.8455; G point outside=False, nonpositive=False, max |line residual|/MC SE=0.6427. national\_\allowbreak{}share65: grid=[0.0, 0.02, 0.04, 0.06, 0.09]; RA point outside=False, nonpositive=False, max |line residual|/MC SE=1.134; G point outside=False, nonpositive=False, max |line residual|/MC SE=0.7098. reporting\_\allowbreak{}shock: grid=[0.0, 0.02, 0.04, 0.06, 0.09]; RA point outside=False, nonpositive=False, max |line residual|/MC SE=0.8369; G point outside=False, nonpositive=False, max |line residual|/MC SE=0.9846. rho\_\allowbreak{}h005: grid=[0.0, 0.02, 0.04, 0.06, 0.09]; RA point outside=False, nonpositive=False, max |line residual|/MC SE=1.481; G point outside=False, nonpositive=False, max |line residual|/MC SE=1.628. rho\_\allowbreak{}h02: grid=[0.0, 0.02, 0.04, 0.06, 0.09]; RA point outside=False, nonpositive=False, max |line residual|/MC SE=1.184; G point outside=False, nonpositive=False, max |line residual|/MC SE=1.261. state\_\allowbreak{}persistence50: grid=[0.0, 0.02, 0.04, 0.06, 0.09]; RA point outside=False, nonpositive=False, max |line residual|/MC SE=3.259; G point outside=False, nonpositive=False, max |line residual|/MC SE=3.007. true2.2\_\allowbreak{}analyst36: grid=[0.0, 0.02, 0.04, 0.06, 0.09]; RA point outside=False, nonpositive=False, max |line residual|/MC SE=1.281; G point outside=False, nonpositive=False, max |line residual|/MC SE=1.023. true2.9\_\allowbreak{}analyst36: grid=[0.0, 0.02, 0.04, 0.06, 0.09]; RA point outside=False, nonpositive=False, max |line residual|/MC SE=0.5604; G point outside=False, nonpositive=False, max |line residual|/MC SE=1.366. true4.3\_\allowbreak{}analyst36: grid=[0.0, 0.02, 0.04, 0.06, 0.09]; RA point outside=False, nonpositive=False, max |line residual|/MC SE=1.162; G point outside=False, nonpositive=False, max |line residual|/MC SE=1.198.
\end{minipage}
\end{table}

\begin{table}[htbp]
\centering
\small
\caption{给定模型下的设计检验大小与功效}\label{tab:design}
\setlength{\tabcolsep}{2pt}
\begin{tabularx}{\linewidth}{@{}>{\raggedright\arraybackslash}p{0.2400\linewidth}>{\raggedright\arraybackslash}p{0.1400\linewidth}>{\raggedright\arraybackslash}p{0.1400\linewidth}>{\raggedright\arraybackslash}p{0.1400\linewidth}>{\raggedright\arraybackslash}p{0.1400\linewidth}>{\raggedright\arraybackslash}p{0.1400\linewidth}@{}}
\toprule
季数 & 计数倍数 & 大小 [95\% CI] & 功效 [95\% CI] & 有效C/N/A & 大小相容 \\
\midrule
4 & 1 & 0.058 [0.03918, 0.08224] & 0.996 [0.9856, 0.9995] & 500/500/500 & 是 \\
4 & 2 & 0.06 [0.04084, 0.08455] & 0.998 [0.9889, 0.9999] & 500/500/500 & 是 \\
4 & 4 & 0.052 [0.03425, 0.07527] & 0.994 [0.9826, 0.9988] & 500/500/500 & 是 \\
8 & 1 & 0.058 [0.03918, 0.08224] & 1 [0.9926, 1] & 500/500/500 & 是 \\
8 & 2 & 0.038 [0.02303, 0.05871] & 1 [0.9926, 1] & 500/500/500 & 是 \\
8 & 4 & 0.082 [0.05949, 0.1096] & 1 [0.9926, 1] & 500/500/500 & 否 \\
12 & 1 & 0.058 [0.03918, 0.08224] & 1 [0.9926, 1] & 500/500/500 & 是 \\
12 & 2 & 0.07 [0.04924, 0.09601] & 1 [0.9926, 1] & 500/500/500 & 是 \\
12 & 4 & 0.054 [0.03588, 0.0776] & 1 [0.9926, 1] & 500/500/500 & 是 \\
20 & 1 & 0.032 [0.0184, 0.05145] & 1 [0.9926, 1] & 500/500/500 & 是 \\
20 & 2 & 0.058 [0.03918, 0.08224] & 1 [0.9926, 1] & 500/500/500 & 是 \\
20 & 4 & 0.056 [0.03753, 0.07992] & 1 [0.9926, 1] & 500/500/500 & 是 \\
\bottomrule
\end{tabularx}
\par\smallskip
\begin{minipage}{\linewidth}\footnotesize\raggedright
条件区分模拟复用观测季节/计数设计，校准与评估样本独立。有效次数依次为校准/零假设/备择假设。末列表示检验大小区间是否包含名义alpha；保留全部单元，包括大小检验不通过的单元。这些模拟不意味着普适季节门槛或普适80\%功效。 alpha=0.05; replicates/stage=500.
\end{minipage}
\end{table}

\begin{table}[htbp]
\centering
\scriptsize
\caption{分季拟合与报告参与率敏感性}\label{tab:secondary}
\setlength{\tabcolsep}{2pt}
\begin{tabularx}{\linewidth}{@{}>{\raggedright\arraybackslash}p{0.2400\linewidth}>{\raggedright\arraybackslash}p{0.1400\linewidth}>{\raggedright\arraybackslash}p{0.1400\linewidth}>{\raggedright\arraybackslash}p{0.1400\linewidth}>{\raggedright\arraybackslash}p{0.1400\linewidth}>{\raggedright\arraybackslash}p{0.1400\linewidth}@{}}
\toprule
统计量/时期 & n & a / 变化 [95\% CI] & b / 原a & b/a / 调整a & 有限/B \\
\midrule
RA / 2022-23 & 365 & 0.01369 [0.0071, 0.02102] & 0.6967 & 50.91 & 1000/1000 \\
RA / 2023-24 & 499 & 0.006174 [0.002667, 0.01542] & 0.7018 & 113.7 & 1000/1000 \\
RA / 2024-25 & 356 & 0.01268 [0.005687, 0.02276] & 0.5826 & 45.93 & 1000/1000 \\
RA / 2025-26 & 528 & 0.02253 [0.008432, 0.03509] & 0.3533 & 15.68 & 1000/1000 \\
RA / 报告调整 & 1748 & 0.0002939 [-0.0003183, 0.0008476] & 0.01495 & 0.01466 & 1000/1000 \\
G / 2022-23 & 367 & 0.07179 [0.03786, 0.1162] & 2.702 & 37.64 & 1000/1000 \\
G / 2023-24 & 503 & 0.03046 [0.0118, 0.07462] & 2.382 & 78.2 & 1000/1000 \\
G / 2024-25 & 357 & 0.06469 [0.03191, 0.1278] & 2.082 & 32.19 & 1000/1000 \\
G / 2025-26 & 529 & 0.1383 [0.03715, 0.2252] & 0.7411 & 5.359 & 1000/1000 \\
G / 报告调整 & 1756 & 0.000373 [-0.003559, 0.003454] & 0.08401 & 0.08364 & 1000/1000 \\
夏季 2023 & 772 & 15.55 & 94.24 & 0.4873 & -- \\
夏季 2024 & 400 & 17.12 & 45.59 & 0.6683 & -- \\
夏季 2025 & 839 & 16.25 & 92.31 & 0.4882 & -- \\
\bottomrule
\end{tabularx}
\par\smallskip
\begin{minipage}{\linewidth}\footnotesize\raggedright
季节行报告a及其区间，b与b/a仅为点估计。报告调整行的a列为原值减调整后a [95\% CI]，b列为原a，b/a列为调整后a；保持原分母不变。夏季行相应三列为计数尺度中位数、报告参与率中位数(\%)及IQR(log G)。唯一自愿报告夏季仅作描述，不作因果或独立重复干预解释。 RA/2022-23: 1000/1000 有限重抽样; RA/2023-24: 1000/1000 有限重抽样; RA/2024-25: 1000/1000 有限重抽样; RA/2025-26: 1000/1000 有限重抽样; G/2022-23: 1000/1000 有限重抽样; G/2023-24: 1000/1000 有限重抽样; G/2024-25: 1000/1000 有限重抽样; G/2025-26: 1000/1000 有限重抽样.
\end{minipage}
\end{table}

